% problem-000267 9 N-酰基亚胺盐反应
\section*{9. N-酰基亚胺盐反应}
N-酰基亚胺盐是一类重要的双亲电试剂，具有丰富的化学反应性质，作为合成子在有机合成中起着重要的作用。

9.1 以下转换展现了很好的反应产率。仔细分析该过程，回答相关问题：

9.1.1 画出中间产物 $\mathbf{A}(\mathrm{C}_{12}\mathrm{H}_{14}\mathrm{ClN}_3\mathrm{O})$ 的结构式。

解答:

（图：）

2 分，其他答案不得分。

9.1.2 画出由 A 转化为最终产物的关键中间体。

解答:

（图：）

每个 2 分，共 8 分；

给分原则：从答案中选择符合要求的给分。骨架错误的不得分。

9.2 如下反应有很好的非对映选择性:

（图：）

9.2.1 最终的关环产物中标黑的两个氢之间的关系是：(a) 顺式；(b) 反式；(c) 无法判断。解答：

a, 1 分，其他答案不得分。

9.2.2 用苯代替甲苯，写出上述反应过程中的关键中间体(要求立体化学)。
解答：

（图：）

（图：）

一共4个，每个2分，共8分；

给分原则：骨架 1 分，立体化学 1 分；骨架正确，立体化学错误得 1 分。立体化学只要求五元环的三个取代基为顺式即可。骨架错误得 0 分。

9.2.3 用末端烯烃替代炔烃，画出此反应产物的结构式(要求立体化学)。

（图：）

$
\mathsf {B F} _ {3} \cdot \mathsf {O E t} _ {2}, \mathrm{甲苯}
$

$

\begin{itemize}
  \item 7 8 ^ {\circ} \mathrm{C}, 5 \mathrm{min}
\end{itemize}

$

解答:

产物结构为:

每个 2 分，共 4 分；如果没有准确标出甲基取代方式，立体化学准确的得 2 分，立体化学不准确继续扣分；产物为以下结构的任意一个，且立体化学准确的可得 1 分

（图：）

（图：）

给分原则：

骨架 1 分，立体化学 1 分；骨架正确，立体化学错误得 1 分，立体化学只要求五元环的三个取代基为顺式即可。不要求甲苯取代基的立体化学。甲基取代基在间位不得分。

9.3 研究结果表明，如下反应有很好的对映选择性。

（图：）

9.3.1 确定此反应原料中 C1 和 C2 的绝对构型。
解答:

C1: 1S, C2: 2R

给分原则：每个1分，共2分，其他答案不得分。

9.3.2 画出中间产物 A 的结构式(要求示出手性中心的绝对构型)。
解答:

or
（图：）

\subsection*{共2分}

没有形成亚胺正离子，形成以下结构的得 1 分

（图：）

\subsection*{给分原则：}

阴离子可以不写，骨架 2 分，苯基和酯基的立体化学错误得 1 分。羰基与酯基必须成反式。骨架错误得 0 分。

9.3.3 判断化合物 B 中 C5 位 H 与酯基的立体关系，并说明理由。
解答：

（图：）

顺式：1 分；

成环的构型取决于酯基；酯基为了避免和酰胺中羰基的排斥，取直立键，或写 1,3-排斥作用；成环时红色的 2 个 H 取反式。2 分。共 3 分

给分原则:

其他答案不得分。

9.4 仔细观察下述反应的底物结构，完成以下反应式，并回答相关问题：

（图：）

9.4.1 画出产物的立体结构式。

解答:

（图：）

（图：）

（图：）

六元环并六元环骨架正确得 1 分，直立键上的 H 与酯基顺式：1 分；苯基与酯基为顺式：1 分，共 3 分；给分原则：

骨架错误，后续正确都不得分。

9.4.2 画出反应的过渡态结构式。

解答:

过渡态结构:

（图：）

骨架正确 1 分；立体化学 1 分，过渡态符号（方括号和≠）1 分；苯基、Cbz、直立键上 H 和甲酯基均处在顺式，如有任何一个立体化学关系错误得 0 分，共 3 分。

\subsection*{给分原则：}

骨架错误，后续正确都不得分。

9.4.3 并依据此过渡态结构分析，与题 9.3 相比，题 9.4 产物中 H 与酯基的立体化学关系是否有变化？(a) 有变化；(c) 没有变化。

解答:

b 没有变，还是 trans 为主；1 分。

9.4.4 此转换的非对映过量值(dr)是升高还是降低？(a) 升高；(b) 降低；(c) 没有变化。
解答：

b，1 分。dr 值降低，

9.4.5 为 9.4.3 的结论给出具体解释。

解答:

酰胺的羰基与酯基成反式，使得亚胺正离子的加成反应从环下方或从酯基的反位进行，与苯基的立体化学无关。共1分。

（图：）

也可以

\subsection*{2025 年 9 月 11 日评分细则}
